\documentclass[10pt,a4paper]{amsart}
\usepackage[margin=19mm]{geometry}
\usepackage{amsmath,amssymb,mathtools}
\usepackage[scr=rsfs,cal=euler]{mathalfa}
\usepackage{xfrac}
\usepackage[unicode,hidelinks,bookmarks=false]{hyperref}
\newcommand{\ie}{i.e.\ }
\newcommand{\cf}{cf.\ }
\newcommand{\resp}{respectively\ }
\newcommand{\Subsection}{\subsection}
\providecommand{\nobreakdash}{\nobreak\hbox{-}}
\setlength{\emergencystretch}{2em}
\multlinegap0pt
\usepackage{amssymb}
\usepackage{tikz}
\usepackage{float}
\newtheorem{proposition}{Proposition}
\def\C {{\mathbb C}}
\title{Minimal resolutions of toric substacks by line bundles}
\author{Zengrui Han}
\date{Source: arXiv:2604.14384v3 (CC BY 4.0)}
\begin{document}
\maketitle

\noindent\emph{Source and licence.} Minimal resolutions of toric substacks by line bundles. Version arXiv:2604.14384v3 (CC BY 4.0), distributed by arXiv under the Creative Commons Attribution 4.0 International licence, http://creativecommons.org/licenses/by/4.0/, as recorded in the arXiv licence metadata for this exact version and shown on the abstract page https://arxiv.org/abs/2604.14384v3. Full licence notice: \texttt{licenses/arxiv-2604.14384-CC-BY-4.0.txt}.\par\smallskip

\noindent\textit{Editorial scope.} Counted unit: the article's proposition prop:MP-inverse-entry (source0.tex lines 441--447). The article's complete original proof of it is delivered with it (source0.tex lines 448--469, one \texttt{proof} environment of 22 lines). No mathematics is written, repaired or re-derived: the statement and the proof are the article's own lines, in the article's own notation, and no retained line is substituted or re-flowed.  No further source-local context is needed: every cross-reference of every retained line resolves inside the two delivered spans.  Nothing was omitted from any retained span, so the package carries no omission record.  No retained line contains a citation, so the wrapper carries no bibliography and no bibliography entry.  No retained line contains a figure environment, an image-inclusion command or drawing code, so no image is delivered and the package declares no asset dependency.  Editorial formatting only, in addition: an AMS \texttt{amsart} preamble that restates the article's own declarations and macros used by the retained lines, namely the environments \texttt{proposition} on separate counters, together with the standard packages those lines need.  The retained environments print in this document as Proposition 1 in order of appearance, whereas the article prints the same items with the numbering of its own full document.  The complete native source of the article as distributed in the arXiv e-print for this exact version is bundled unchanged as \texttt{sources/198534/source0.tex}.
\begin{proposition}\label{prop:MP-inverse-entry}
    Let $\partial:\C^n\rightarrow\C^m$ be a linear map and $\partial^+$ be its Moore-Penrose inverse. Then the entry of $\partial^+$ at the $i$-th row and $j$-th column is given by the following formula
    \begin{align*}
        \partial^+_{ij} = \frac{1}{\Delta}\sum_{\substack{ ST \text{ hedge} \\ j\in S, i\in T}} (-1)^{a+b} \overline{\det(\partial_{S\times T})}\det(\partial_{S\backslash\{j\}\times T\backslash\{i\}}).
    \end{align*}
    where $a$ and $b$ are the positions of $i$ in $T$ and $j$ in $S$ respectively.
\end{proposition}

\begin{proof}
    Denote the standard basis vectors of $\C^n$ and $\C^m$ by $\{e_i\}$ and $\{f_j\}$ respectively. Then the entry $\partial^+_{ij}$ is given by
    \begin{align*}
        \partial^+_{ij} = \langle e_i, \partial^+(f_j) \rangle = \frac{1}{\Delta}\sum_{ST}|\det\partial_{S\times T}|^2 \langle e_i, \partial_{ST}^+(f_j) \rangle
    \end{align*}
    according to the hedge formula. Now we claim
    \begin{align*}
        \langle e_i, \partial_{ST}^+(f_j) \rangle = (\partial_{S\times T}^{-1})_{\mathrm{pos}_T(i),\mathrm{pos}_S(j)}
    \end{align*}
    if $j\in S$ and $i\in T$ and is zero otherwise, where $\mathrm{pos}_T(i)$ and $\mathrm{pos}_S(j)$ are the positions of $i$ in $T$ and $j$ in $S$ respectively. To see this, first note that if $j\not\in S$ then by definition $\partial_{ST}^+(f_j)=0$, hence we assume $j\in S$. Since $\partial_{S\times T}:\C\{T\}\rightarrow\C\{S\}$ is an isomorphism, there exists a unique $x_j\in \C\{T\}$ such that $\partial_{S\times T}(x_j)=f_j$. In matrix form, we have
    \begin{align*}
        x_j = \sum_{i\in T}(\partial_{S\times T}^{-1})_{\mathrm{pos}_T(i),\mathrm{pos}_S(j)} e_i.
    \end{align*}
    Consider $f_j = (f_j - \partial x_j) + \partial x_j$. By construction the first term is in $\C\{\overline{S}\}$, and the second in the span of $\partial(T)$. Applying $\partial_{ST}^{+}$ we get
    \begin{align*}
        \partial_{ST}^{+}(f_j) = 0 + \partial_{ST}^{+}(\partial x_j)= x_j = \sum_{i\in T}(\partial_{S\times T}^{-1})_{\mathrm{pos}_T(i),\mathrm{pos}_S(j)} e_i
    \end{align*}
    and the claim follows. The proposition now follows from the adjugate matrix formula
    \begin{align*}
        (\partial_{S\times T}^{-1})_{\mathrm{pos}_T(i),\mathrm{pos}_S(j)} = (-1)^{a+b} \frac{\det\left( \partial_{S\backslash\{j\} \times T\backslash\{i\}} \right)}{\det\left( \partial_{S\times T} \right)}.
    \end{align*}
\end{proof}

\end{document}
